%%====================================================================
%% Single source for manuscript_clean.tex and manuscript_marked.tex.
%% Blocks between %%REV-START and %%REV-END are new or substantively
%% revised in this round; build_manuscripts.py strips the markers for the
%% clean file and colours the blocks red for the marked file.
%%====================================================================
\documentclass[fleqn,10pt]{wlscirep}
\usepackage[utf8]{inputenc}
\usepackage[T1]{fontenc}
\usepackage{siunitx}
\usepackage{placeins}
\graphicspath{{figures/}}
%%REV-START
\title{BioInteract: Interpretable Drug--Target Interaction Prediction via
  Residue-Level Cross-Attention with Protein-Language and Physicochemical Priors}
%%REV-END

\author[1,+]{Shengnan Wang}
\author[2,+]{Qi Zhang}
\author[2]{Xucheng Liu}
\author[2]{Zixuan Wang}
\author[2]{Jiangke Cheng}
\author[3,4]{Zhi Huang*}
\affil[1]{School of Biological and Chemical Engineering,
  Panzhihua University, Panzhihua 617000, P.R.\ China}
\affil[2]{School of Mathematics and Computer Science,
  Panzhihua University, Panzhihua 617000, P.R.\ China}
\affil[3]{Department of Environmental Science and Engineering,
  College of Architecture and Environment,
  Sichuan University, Chengdu 610065, P.R.\ China}
\affil[4]{Institute of New Energy and Low Carbon Technology,
  Sichuan University, Chengdu 610065, P.R.\ China}

\affil[*]{corresponding.author. E-mail: 525514353@qq.com}

\affil[+]{these authors contributed equally to this work}

\begin{abstract}
%%REV-START
Accurate prediction of drug--target interactions (DTIs) can support the
prioritisation of candidate pairs, yet published models are usually compared
through numbers quoted across papers whose protocols differ, and their
attribution outputs are rarely accompanied by an explicit statement of where
the model may be trusted. We present BioInteract, a DTI framework combining a
chemistry-aware molecular graph encoder, an ESM-2/physicochemical protein
residue encoder with a configured domain-label channel, and bidirectional
cross-attention, and we evaluate it inside a matched-protocol benchmark on the
Davis kinase inhibitor dataset in which every model is trained with one shared
partition, seed, optimiser, schedule and threshold rule. The comparison is a
three-model ladder that changes one component at a time, so that adding the
frozen ESM-2 representation and then replacing concatenation with
cross-attention can be read separately. Both steps help, but by
protocol-dependent amounts: the protein language model adds 0.120 and 0.063
AUPRC under the random and Target-ID-held-out protocols, whereas cross-attention
adds 0.067 and 0.201 AUPRC under the same two protocols. A breadth check across two
further architecture families places these differences in context: on the random
split the five architectures tested fall within a narrow band, so the
entity-held-out protocols carry the informative contrasts. We audit chemical and
sequence similarity across all partitions, evaluate exact-sequence-grouped
cold-start protocols, and report leakage-aware diagnostic controls. On a strict
exact-entity-novel BindingDB cohort, whose construction is documented as a
complete row-level filtering waterfall for both the curated-articles snapshot
and the full BindingDB release, the frozen Davis checkpoint achieved an AUROC of
0.560 and an AUPRC of 0.340; an applicability-domain analysis attributes this
limited transfer principally to ligand and target novelty relative to Davis.
The attribution analyses generate hypotheses about model-relevant atoms and
residues, but, in line with evidence from models trained with explicit
non-covalent interaction supervision, unsupervised attention is not treated
here as a recovered interaction mechanism. BioInteract therefore contributes an
audited, reproducible and domain-of-applicability-aware DTI reference rather
than a new architectural principle.
%%REV-END
\end{abstract}
\begin{document}

\flushbottom
\maketitle
\thispagestyle{empty}
\section*{Keywords}

Drug--target interaction; interpretable deep learning; cross-attention; graph
neural network; protein language model; model attribution
%%====================================================================
\section{Introduction}
\label{sec:introduction}
%%====================================================================

The emergence of ``AI for Science'' has increased expectations that
computational models should not merely maximise predictive accuracy, but should
also expose the evidence and uncertainty needed for scientific hypothesis
generation.\cite{schneider2020rethinking}  This expectation is especially
important in drug discovery, where identifying drug--target interactions (DTIs)
requires reliable prioritisation together with inspectable outputs that can
guide, but not replace, mechanistic follow-up.

Experimental profiling of binding affinities through surface plasmon resonance,
isothermal titration calorimetry, or kinase activity assays is laborious and
prohibitively expensive.  The Davis kinase inhibitor
dataset\cite{davis2011comprehensive} alone required systematic measurement
of 72 inhibitors against 442 kinases---a campaign spanning only a fraction of
the druggable proteome.  Computational DTI prediction can ease this bottleneck
by prioritising the most promising candidates and thereby reducing cost and
attrition in development
pipelines.\cite{chen2016drug,yamanishi2008prediction}  Yet the field has
largely remained trapped in a \emph{predictive-accuracy race}: successive deep
learning architectures---DeepDTA,\cite{ozturk2018deepdta}
GraphDTA,\cite{nguyen2021graphdta}
MolTrans,\cite{huang2021moltrans}
DrugBAN\cite{bai2023drugban}---have incrementally improved benchmark metrics
while leaving an important modelling question only partly addressed:
\emph{which input features drive a given prediction?}

%%REV-START
Pairwise DTI modelling is only one of several complementary routes to the same
practical decision.  Ligand-based affinity and selectivity profiling can
prioritise modulators of a receptor family while explicitly delimiting the
applicability domain of the resulting models and validating them on external
compounds, as illustrated by ALPACA for cannabinoid receptor
modulators;\cite{delre2023alpaca} unified frameworks such as DTIAM additionally
couple interaction prediction with affinity regression and mechanism-of-action
assignment.\cite{lu2025dtiam}  These examples are cited here for context: they
indicate that applicability-domain assessment and external validation, rather
than benchmark ranking alone, determine whether a predictive model is usable,
and the present study adopts that emphasis for a pairwise DTI classifier.
%%REV-END

This question lies at the heart of the broader ``AI for Science'' agenda.  A
prediction with no interpretable output is, at best, a hypothesis generator.
Moving from \emph{prediction} to hypothesis generation therefore benefits from
models that expose internal attribution patterns---for example, attention
distributed across drug atoms and protein residues and gradient-based atom
scores---rather than returning only a scalar probability.  These outputs can
prioritise follow-up experiments, but they are not direct measurements of
physical contacts or pocket geometry.

Classical DTI methods illuminate the fundamental tension between interpretability
and accuracy.  Structure-based approaches (molecular docking, molecular dynamics)
provide structural inference anchored in three-dimensional protein
geometry,\cite{kitchen2004docking} yet are constrained by the availability of
high-quality crystal structures and the steep computational cost of free-energy
calculations.  Ligand-based methods (QSAR, molecular
fingerprints)\cite{cherkasov2014qsar} circumvent structural requirements at
the expense of mechanistic transparency.  Neither paradigm fully satisfies the
twin demands of modern drug discovery.

Deep learning has transformed predictive capabilities in this domain.
Pretrained protein language models such as ESM-2\cite{lin2023evolutionary}
encode evolutionary and structural information at residue resolution, while
graph neural networks faithfully capture molecular topology.  Concurrently,
chemical language models have demonstrated that pretrained molecular
representations support accurate and interpretable prediction of reactivity and
toxicity,\cite{huang2025large} and chemistry-enhanced transformer
architectures such as Radical-Net recover atom-level reaction
mechanisms.\cite{huang2025radical}  Together, these advances underscore
the importance of domain-specific pretraining and chemistry-aware inductive
biases in molecular AI.\cite{huang2024ai}

In this study, ``prior knowledge'' refers specifically to the
sequence information encoded by frozen ESM-2 pretraining together with explicit
per-residue physicochemical descriptors. It does not denote curated Pfam,
InterPro, binding-site, or other protein-domain annotations, and the manuscript
title has been narrowed accordingly.

The ``AI for Science'' paradigm benefits from attribution mechanisms that are
part of the prediction pathway rather than appended only after training.
Attention visualisation and SHapley Additive exPlanations (SHAP) can nevertheless
have faithfulness limitations.\cite{jain2019attention} A central principle is
therefore that, when attention is an integral component of the prediction
pathway, the resulting attributions can be examined consistently with model
output but still require external validation before mechanistic interpretation.

%%REV-START
\subsection*{Scope of the present contribution}

The individual components used here are established.  Chemistry-aware molecular
graph encoders, residue-level protein language-model features, and cross-modal
attention have each been explored extensively, and combinations of them have
been published for DTI and drug--target affinity (DTA)
prediction.\cite{nguyen2022gefa,bal2023pgraphdta,li2025hcafdta}  We therefore do
not claim that any single element of BioInteract, or their combination, is
methodologically new.  Nor do we claim that its attention maps recover binding
mechanisms: as discussed in Section~\ref{sec:related}, prior work using explicit
non-covalent interaction supervision indicates that attention learned only from
interaction labels need not be biologically
meaningful.\cite{li2020monn,gim2023arkdta,seo2024blendnet}

What this study contributes is an audited reference point for the Davis binary
high-affinity classification task, comprising four elements.  First, a
matched-protocol benchmark (Section~\ref{sec:matched}) in which BioInteract is
compared against re-implemented baselines trained under a single shared
partition, seed, optimiser, schedule, loss and threshold rule.  Its primary form
is a three-model ladder that changes one component at a time---molecular graph
without a protein language model, then with one, then with cross-attention in
place of concatenation---so that each difference has a single cause rather than
reflecting a difference in training recipe or reporting convention.  Second, a leakage audit of the identifier-based
splits together with exact-sequence-grouped cold-start protocols and
leakage-aware diagnostic controls.  Third, a strict external evaluation on
BindingDB whose construction is reported as a complete row-level filtering
waterfall for both the curated-articles snapshot and the full database release,
accompanied by an applicability-domain analysis that quantifies why transfer is
limited.  Fourth, model-native attribution outputs released with an explicit
statement of what they do and do not establish.  Several of these analyses
reduce, rather than support, the apparent strength of the model; we report them
because a reference point is only useful if its boundaries are visible.
%%REV-END

BioInteract rests on three architectural components:

\begin{enumerate}[leftmargin=*]
\item \textbf{Chemistry-aware molecular graph encoding.}  Drug molecules are
  represented as attributed graphs encoding functional-group-relevant properties
  (H-bond donor/acceptor status, hydrophobicity, partial charges) that
  are available to the learned molecular representation.

\item \textbf{Protein residue encoding with a disclosed domain-label limit.}
  Protein representations fuse ESM-2 residue embeddings with per-residue
  physicochemical properties and a configured domain-label channel. No curated
  per-protein domain-annotation files are present in the released Davis input
  package; the channel therefore receives the unknown label at every position
  and is not interpreted as Pfam/InterPro information.

\item \textbf{Bidirectional cross-attention attribution module.}  A multi-head
  cross-attention mechanism computes a matrix
  $\mathbf{M} \in \mathbb{R}^{n_\text{atoms} \times n_\text{residues}}$ of
  learned attention weights between drug-atom and protein-residue
  representations.  The matrix is an integral component of the forward pass;
  it is used here as a model-native attribution output, not as a predicted
  atom-level structural contact map.
\end{enumerate}

We evaluate BioInteract on the Davis kinase inhibitor
dataset\cite{davis2011comprehensive} under a random pair split and two
entity-held-out protocols: Drug-ID-held-out and Target-ID-held-out. We report
chemical and sequence-similarity
diagnostics for these partitions, because identifier-level exclusion alone does
not guarantee scaffold- or sequence-level novelty.  We also add
exact-sequence-grouped cold-target and exact-sequence-grouped cold-both
evaluations.  The attribution
analyses---cross-attention maps, Grad-CAM scores, attention-concentration
summaries, and case studies---are interpreted as model behaviour rather than as
independent evidence of molecular recognition mechanisms.

%%====================================================================
\section{Related Work}
\label{sec:related}
%%====================================================================

\subsection{Deep Learning for DTI Prediction}

Deep learning approaches to DTI prediction have evolved through several
generations.\cite{wen2017deep}  \emph{Sequence-based} models encode both
drug SMILES and protein amino acid sequences as one-dimensional strings.
DeepDTA\cite{ozturk2018deepdta} employs parallel one-dimensional convolutional
neural networks (CNNs) for drug and
protein encoding; WideDTA\cite{ozturk2019widedta} extends this framework
with additional molecular descriptors.  Although computationally efficient, these
models discard molecular topology and are limited in their capacity to model
long-range intramolecular interactions.

\emph{Graph-based} models overcome this limitation by representing drugs as
attributed molecular graphs.  GraphDTA\cite{nguyen2021graphdta} benchmarks a
graph convolutional network (GCN), graph attention network (GAT), graph
isomorphism network (GIN), and hybrid GAT--GCN encoder while retaining a CNN
for proteins.  MolTrans\cite{huang2021moltrans} introduces mutual
interaction learning between molecular substructure representations and protein
subsequences.

\emph{Attention-based} models include architectures that learn local or
cross-modal representations for DTI/DTA prediction.
TransformerCPI\cite{chen2020transformercpi} applies transformer attention to
drug--protein interaction patterns.  DrugBAN\cite{bai2023drugban} introduces
a dual bilinear attention network for fine-grained local interaction modelling.
AttentionDTA\cite{zhao2022attentiondta} employs multi-head self-attention for
both drug and protein representations.
%%REV-START
Because these models were published under different splits, binarisation rules,
thresholds and reporting conventions, values quoted across their papers are not
directly comparable.  Section~\ref{sec:matched} therefore re-implements
representatives of the sequence-CNN, molecular-graph and bilinear-attention
families and trains them under one shared protocol.
%%REV-END

\subsection{Pretrained Protein Language Models}

ESM-2,\cite{lin2023evolutionary} trained on millions of protein sequences
via masked language modelling, produces residue-level embeddings that encode
evolutionary conservation, secondary structure propensity, and contact
information.  These pretrained representations substantially outperform
hand-crafted features on downstream tasks.\cite{rives2021biological}  In DTI
prediction, pretrained representations have already been combined with graph
drug encoders and interaction modules.  For example, GEFA uses pretrained
protein representations in an early-fusion graph architecture,\cite{nguyen2022gefa}
PGraphDTA replaces a CNN target encoder with protein language-model features,
\cite{bal2023pgraphdta} and HCAF-DTA combines protein representations with
bidirectional cross-attention.\cite{li2025hcafdta}  BioInteract therefore does
not claim to be the first application of a protein language model, a graph
encoder, or cross-attention to DTI/DTA prediction.
%%REV-START
Because the effect of protein language-model pretraining cannot be separated
from architectural choices by comparing against models that lack it, the
matched-protocol benchmark in Section~\ref{sec:matched} includes three
baselines that consume the identical frozen \texttt{esm2\_t30\_150M\_UR50D}
representation used by BioInteract and differ only in how drug and protein
representations are combined.
%%REV-END

Recent DTA studies also illustrate the use of residue-level protein graphs,
domain-specific representations, or biological-context features alongside
evaluation on more than one benchmark, including 3DProtDTA,\cite{voitsitskyi2023_3dprotdta}
the domain-specific framework of Liu \textit{et al.},\cite{liu2024domainspecific}
and InceptionDTA.\cite{kalemati2025inceptiondta}
%%REV-START
These studies, while not directly comparable to the present Davis binary
classification protocol, provide further context for alternative approaches to
the DTA problem.
%%REV-END

\subsection{Interpretability in DTI Models}

Existing interpretability approaches include post-hoc methods (attention-weight
visualisation, gradient-based attribution, and SHAP values\cite{jimenez2020drug})
and model-native attribution mechanisms.  Attention weights do not necessarily
equal feature importance.\cite{jain2019attention}
%%REV-START
This general caution has a specific and well-documented form in DTI modelling.
MONN showed that jointly supervising non-covalent interaction labels derived
from structural complexes is needed to obtain pairwise interaction predictions
that agree with those complexes;\cite{li2020monn} ArkDTA reached the same
conclusion by regularising protein--ligand attention towards non-covalent
interaction annotations;\cite{gim2023arkdta} and BlendNet reported that
interaction-level supervision is what allows a complex-structure-free model to
produce attributions consistent with known binding
behaviour.\cite{seo2024blendnet}  We do not treat these studies as a separate
category of work that lies outside the scope of comparison.  We read them as
direct evidence about conventional attention-based models, including
BioInteract: absent explicit interaction supervision, there is no reason to
expect attention weights to identify the residues that actually contact a
ligand, and a small number of qualitative visualisations cannot establish that
they do.  BioInteract accordingly uses its cross-attention matrix as a
model-native attribution that contributes to the prediction and can be inspected
alongside it, and we make no claim that this attribution recovers atom-level
contacts, binding modes, or physical interaction types.  Section~\ref{sec:interpretability}
reports the attribution outputs as descriptions of model behaviour, and
Section~\ref{sec:discussion} states the structural evaluation that would be
required before any mechanistic reading.
%%REV-END

\subsection{Cold-Start Evaluation}

Random-split evaluation permits information leakage through shared drugs or
targets between training and test sets, yielding inflated
estimates.\cite{huang2021therapeutics}  Cold-start protocols address this by
ensuring that test identifiers are absent from training data. We report the
original Drug-ID-held-out and Target-ID-held-out protocols alongside the random
split, and separately use exact-sequence grouping for stricter within-Davis
stress tests.

%%====================================================================
\section{Methods}
\label{sec:methods}
%%====================================================================

\subsection{Problem Formulation}

Given a drug molecule $d$ represented by its SMILES string and a target protein
$t$ represented by its amino acid sequence, we formulate DTI prediction as
binary classification: predict whether $d$ binds $t$ with high affinity.
Following the established protocol for the Davis
dataset,\cite{davis2011comprehensive} continuous $K_d$ values are binarised
at \SI{30}{nM} ($pK_d > 7.52$), classifying pairs with $K_d < \SI{30}{nM}$ as
positive interactions.

\subsection{Architecture Overview}

BioInteract comprises five sequential modules:
(1)~a graph isomorphism network with edge features (GINE)-based drug encoder producing per-atom representations,
(2)~a target encoder combining ESM-2 embeddings with physicochemical features
and a configured domain-label channel,
(3)~a bidirectional cross-attention module computing attention matrices between
ligand-atom and protein-residue representations,
(4)~gated attention pooling to aggregate variable-length representations, and
(5)~a multilayer perceptron (MLP) prediction head.
The full model contains 2{,}442{,}083 trainable parameters.

\subsection{Drug Molecular Graph Encoder}

\subsubsection{Chemistry-Aware Atom Featurisation}

Each atom is represented by a 52-dimensional feature vector encoding:
\begin{itemize}[leftmargin=*]
\item \textbf{Topological features} (38 dimensions): atom type (one-hot,
  20 types), hybridisation state (6), formal charge (6), number of
  attached hydrogens (6).
\item \textbf{Structural features} (9 dimensions): degree, aromaticity, ring
  membership, and ring size membership (3--8).
\item \textbf{Functional features} (5 dimensions): hydrogen-bond donor and
  acceptor status, Crippen logP contribution, molar refractivity, and Gasteiger
  partial charge.
\end{itemize}
Bond features (16 dimensions) encode bond type, conjugation, ring membership,
and stereochemistry.

\subsubsection{GINE Message Passing}

We employ a 3-layer GINE
architecture:\cite{hu2020strategies}
\begin{equation}
\widetilde{\mathbf{e}}_{vu}=\mathbf{W}_e\mathbf{e}_{vu}+\mathbf{b}_e
\in\mathbb{R}^{256},\qquad \mathbf{e}_{vu}\in\mathbb{R}^{16},
\label{eq:edge_projection}
\end{equation}
\begin{equation}
\mathbf{h}_v^{(\ell+1)} = \mathrm{MLP}^{(\ell)}\!\left(
  \mathbf{h}_v^{(\ell)} +
  \sum_{u\in\mathcal{N}(v)}
  \mathrm{ReLU}\!\left(\mathbf{h}_u^{(\ell)} + \widetilde{\mathbf{e}}_{vu}\right)
\right)
\end{equation}
where $\mathbf{h}_v^{(\ell)}$ is the representation of atom $v$ at layer $\ell$,
$\mathbf{e}_{vu}$ denotes the raw 16-dimensional edge feature vector, and
$\widetilde{\mathbf{e}}_{vu}$ is its learned 16-to-256 projection before message
passing. The released \texttt{GINEConv} instantiation uses the default fixed
$\epsilon=0$. Each layer incorporates batch normalisation, ReLU activation,
dropout ($p=0.2$), and a residual connection; the hidden dimension is $D=256$.
Global graph pooling is \emph{not} applied at this stage, preserving individual
atom representations for the subsequent cross-attention step.

The archived checkpoints use the stated molecular featurisation; public
checkpoint reconstruction applies no stochastic graph augmentation.

\subsection{Protein Target Encoder}

\subsubsection{ESM-2 Residue Embeddings}

Residue-level representations are extracted from ESM-2 (150\,M parameters,
\texttt{esm2\_t30\_150M\_UR50D}),\cite{lin2023evolutionary} yielding
640-dimensional embeddings for each amino acid.  These embeddings encode
evolutionary conservation, secondary structure propensity, and long-range contact
information learned from 65 million protein sequences.  ESM-2 parameters are
frozen throughout training.

\subsubsection{Physicochemical Property Encoding}

Each residue is augmented with a 4-dimensional physicochemical vector:
\begin{equation}
\mathbf{p}_i = \bigl[\text{hydrophobicity}_i,\;\text{mol.\ weight}_i,\;
  \text{charge}_i,\;\text{polarity}_i\bigr]
\end{equation}
Hydrophobicity follows the Kyte--Doolittle scale.\cite{kyte1982simple}
These values provide residue-level physicochemical descriptors; by themselves,
they do not identify pocket membership, solvent exposure, or interaction
partners.

\subsubsection{Configured Domain-Label Channel and Availability}

The implementation contains a learnable domain-label embedding layer:
\begin{equation}
\mathbf{d}_i = \mathrm{DomainEmbed}(\text{domain\_type}_i) \in \mathbb{R}^{32}
\end{equation}
The label builder assigns the unknown \texttt{NONE} type when a per-protein
annotation file is unavailable. The released Davis package contains no such
annotation files, so all archived input positions use this unknown label. We
therefore do not represent this channel as curated Pfam/InterPro context, as a
binding-site label, or as evidence for a domain-specific conclusion.

\subsubsection{Feature Fusion}

The three feature sources are concatenated and projected through a shared linear
transformation:
\begin{equation}
\mathbf{r}_i = \mathrm{Proj}\!\left(
  [\mathbf{e}_i \,\|\, \mathbf{p}_i \,\|\, \mathbf{d}_i]
\right) \in \mathbb{R}^{D}
\end{equation}
where $\mathbf{e}_i \in \mathbb{R}^{640}$ is the ESM-2 embedding and Proj is a
two-layer MLP with layer normalisation and Gaussian error linear unit (GELU)
activation. In the released configuration, $\mathbf{d}_i$ is the shared unknown
domain-label embedding rather than a curated per-residue annotation. The projected
dimension $D=256$ matches the drug encoder output.

\subsection{Cross-Attention Interaction Module}

The cross-attention module produces the model-native attribution used in this
study. It computes bidirectional attention between ligand-atom and
protein-residue representations, producing an explicit atom--residue attention
matrix.

\subsubsection{Drug-to-Protein Attention}

\begin{equation}
\mathrm{Attn}(\mathbf{Q}_d,\mathbf{K}_p,\mathbf{V}_p) =
  \mathrm{softmax}\!\left(\frac{\mathbf{Q}_d\mathbf{K}_p^\top}{\sqrt{d_k}}\right)
  \mathbf{V}_p
\end{equation}
where $\mathbf{Q}_d = \mathbf{H}_d\mathbf{W}_Q^d$,
$\mathbf{K}_p = \mathbf{R}\mathbf{W}_K^p$,
$\mathbf{V}_p = \mathbf{R}\mathbf{W}_V^p$, with $H=8$ heads and head
dimension $d_k = D/H = 32$.

\subsubsection{Protein-to-Drug Attention}

Symmetrically, protein residues attend over drug atoms:
\begin{equation}
\mathrm{Attn}(\mathbf{Q}_p,\mathbf{K}_d,\mathbf{V}_d) =
  \mathrm{softmax}\!\left(\frac{\mathbf{Q}_p\mathbf{K}_d^\top}{\sqrt{d_k}}\right)
  \mathbf{V}_d
\end{equation}

\subsubsection{Attribution Map Extraction}

The atom--residue attribution map is obtained by averaging drug-to-protein
attention weights across all heads:
\begin{equation}
\mathbf{M} = \frac{1}{H}\sum_{h=1}^{H}
  \mathrm{Attn}_h(\mathbf{Q}_d,\mathbf{K}_p)
  \;\in\;\mathbb{R}^{n\times L}
\label{eq:interaction_map}
\end{equation}
Because $\mathbf{M}$ is computed within the forward pass rather than post-hoc,
it is a model-native attribution matrix.  It is not interpreted as a physical
contact map or as an assignment of a specific interaction type to an atom pair.

\subsection{Gated Pooling and Prediction}

Variable-length representations are aggregated via gated attention pooling:
\begin{equation}
\mathbf{z} = \sum_{i}\alpha_i\cdot g(\mathbf{x}_i),\qquad
\alpha_i = \frac{\exp(f(\mathbf{x}_i))}{\sum_j\exp(f(\mathbf{x}_j))}
\end{equation}
where $f\colon\mathbb{R}^D\to\mathbb{R}$ and
$g\colon\mathbb{R}^D\to\mathbb{R}^D$ are learned two-layer MLPs with Tanh
activation.  The pooled drug and protein representations are concatenated and
passed through a two-hidden-layer MLP (256\,$\to$\,128; ReLU; dropout $p=0.3$)
followed by sigmoid activation to produce the binding probability.

%%REV-START
\subsection{Matched-Protocol Comparison Models}
\label{sec:matched_methods}

To compare architectures without importing values measured under other
protocols, we re-implemented a set of baselines and trained each of them,
together with BioInteract, inside one harness.  Every run uses the same Davis records,
the same partition for a given protocol, seed 42, weighted binary cross-entropy
with 0.05 label smoothing, AdamW (learning rate $10^{-4}$, weight decay
$10^{-5}$), an effective batch size of 64 formed from micro-batches of 16 with
gradient accumulation, automatic mixed precision, gradient-norm clipping at 1.0,
a five-epoch linear warm-up followed by cosine decay, a maximum of 30 epochs
with early stopping at patience 6 on validation AUROC, and a decision threshold
selected once on that protocol's validation predictions.  Cross-attention
models allocate one atom--residue map per pair, so the micro-batch and
mixed-precision settings bound accelerator memory; validation and test
inference are run in full precision so that reported metrics do not depend on
the mixed-precision training path.

The primary comparison is a three-model ladder in which exactly one component
changes at each step, so that each difference has a single cause:

\begin{itemize}[leftmargin=*]
\item \textbf{GraphDTA}\cite{nguyen2021graphdta} --- the molecular graph
  encoder with a character convolutional protein encoder.  Molecular graph but
  no protein language model.
\item \textbf{ESM2-GraphConcat} --- the same GINE atom encoder, with the
  convolutional protein encoder replaced by BioInteract's
  ESM-2/physicochemical residue encoder; the two representations are pooled
  independently and concatenated.  It differs from GraphDTA only by the protein
  language model.
\item \textbf{BioInteract} --- identical encoders to ESM2-GraphConcat, with
  concatenation replaced by bidirectional cross-attention.  It differs from
  ESM2-GraphConcat only by the interaction module.
\end{itemize}

The second and third rows share BioInteract's frozen protein representation
exactly.  The first-to-second step therefore isolates the contribution of ESM-2
pretraining under a fixed training recipe and a fixed drug encoder, and the
second-to-third step isolates the contribution of the cross-modal interaction
module.  Because GraphDTA is a re-implementation from its published
architecture description rather than a run of the authors' released code, it
should be read as a controlled reference for its architecture family rather
than as a reproduction of the specific numbers reported in that paper.  Every
comparison reported in this manuscript was measured in this study under the
protocol described above; no value is carried over from another publication.

Two further architectures---a sequence convolutional model in the DeepDTA
family\cite{ozturk2018deepdta} and the low-rank bilinear attention and bilinear
pooling introduced by DrugBAN\cite{bai2023drugban} applied to the same two
encoders---were trained by the same harness on the random split as a breadth
check across model families.  They are reported in Supporting Information
Table~S14 and are discussed in Section~\ref{sec:matched}.
%%REV-END

\subsection{Training Provenance and Evaluation Thresholds}

The archived entity-split checkpoints preserve model architecture but do not
store a uniform optimizer, scheduler, warmup, or label-smoothing record. The
Random and Drug-ID-held-out logs explicitly record weighted binary
cross-entropy, label smoothing of 0.05, AdamW, and a final-run header, whereas
the checkpoint-matched Target-ID-held-out trace has the older
\texttt{run\_split.py} header and records only loss and validation AUROC. Its
checkpoint does not retain optimizer or loss state. We therefore do not
retrospectively assert a single training recipe across all three historical
checkpoints (Supporting Information, Training-provenance table).

No historical entity-split model is retrained for the current-release metrics.
The canonical reconstruction instantiates each of those three released
checkpoints from its embedded model configuration, performs full-precision
inference, selects an F1 threshold once on the corresponding validation
predictions, and applies that frozen threshold to test predictions. This
separation prevents test-set threshold selection.
%%REV-START
These archived checkpoints are the source of the ablation results in
Section~\ref{sec:ablation} only.  All architecture comparisons in
Section~\ref{sec:matched} come from the matched-protocol harness described
above and are independent of the archived training histories.
%%REV-END

In contrast, the two exact-sequence-grouped strict protocols were trained from
initialisation on their respective new training partitions. They use seed 42,
the same Davis preprocessing and architecture, weighted binary cross-entropy
with 0.05 label smoothing, AdamW (learning rate $10^{-4}$; weight decay
$10^{-5}$), batch size 64, gradient-norm clipping at 1.0, a five-epoch linear
warm-up followed by cosine decay, and a maximum of 100 epochs. The checkpoint
with the highest validation AUROC was selected with patience 15, after which an
F1 threshold was selected on that protocol's validation predictions and frozen
for its test evaluation. Complete strict-split training provenance, checkpoint
paths, SHA-256 hashes, selected epochs, and thresholds are provided in the
Supporting Information Training-provenance table.

\subsection{Interpretability Analysis}

\subsubsection{Residue-Level Attention Profiling}

For a given drug--target pair, per-residue attention scores are obtained by
summing the interaction map over all drug atoms:
\begin{equation}
s_j = \sum_{i=1}^{n} M_{ij},\quad j=1,\ldots,L
\end{equation}
These values are model-attribution scores and are not assigned a biological
hotspot label.

\subsubsection{Grad-CAM Functional-Group Attribution}

Grad-CAM\cite{selvaraju2017grad} is adapted to the graph-neural-network (GNN)
setting:
\begin{equation}
\mathrm{CAM}_v = \mathrm{ReLU}\!\left(\sum_k w_k\cdot A_v^k\right),\quad
w_k = \frac{1}{|\mathcal{V}|}\sum_v\frac{\partial z}{\partial A_v^k}
\end{equation}
where $A_v^k$ is channel $k$ of the third (final) GINE message-passing layer
and $z$ is the pre-sigmoid positive-class logit. Gradients are therefore taken
with respect to the logit, not the post-sigmoid probability.
Functional groups are detected with predefined SMARTS patterns (aromatic ring,
hydroxyl, amino, amide, carboxyl, sulfonyl, halogen, ether, heterocycle N, and
carbonyl). For each group, all RDKit substructure matches are merged into a
deduplicated atom set and the group score is the mean of the normalised
atom-level scores over that set. An atom can match more than one group pattern.
These scores indicate sensitivity of the prediction to the learned
representation and do not identify chemical bonds or contacts in a
three-dimensional complex.

\subsubsection{Attention Sparsity}

For each drug--target pair, residue scores are normalised by that pair's maximum:
\begin{equation}
\widetilde{s}_j = \frac{s_j}{\max_k s_k}.
\end{equation}
Pair-level sparsity is then
\begin{equation}
\mathrm{Sparsity} = \frac{|\{j : \widetilde{s}_j < 0.1\}|}{L}.
\end{equation}
The reported global value uses all 1{,}506 positive Davis records, rather than a
single test partition, evaluated with the archived \texttt{checkpoints/best.pt}
checkpoint. Each pair is normalised separately; the 1{,}301{,}380 valid-residue
scores are then concatenated before the global distribution and fraction are
computed. Consequently, all original train/validation/test memberships are
represented and longer valid target sequences contribute more residue-level
observations.

%%====================================================================
\section{Experimental Setup}
\label{sec:experiments}
%%====================================================================

\subsection{Dataset}

The Davis kinase inhibitor
dataset\cite{davis2011comprehensive} comprises 68 drugs (kinase inhibitors)
and 442 targets (protein kinases) forming 30{,}056 drug--target pairs.  Binding
affinities ($K_d$) are binarised at \SI{30}{nM}, yielding approximately
\SI{5}{\percent} positive interactions.

\subsection{Pre-specified External BindingDB Evaluation}
\label{sec:bindingdb_methods}

To test exact-entity transfer without changing the submitted Davis training
protocol, we curated the official BindingDB curated-articles snapshot
\texttt{BindingDB\_BindingDB\_Articles\_202607}.\cite{gilson2016bindingdb}
%%REV-START
This snapshot is the subset of BindingDB curated in-house from the primary
literature, and it contains 93{,}712 measurement rows rather than the millions
of rows in the complete release.  Table~\ref{tab:bindingdb_curation} reports the
full row-level accounting for both this snapshot and, as an independent
verification requested during review, the complete
\texttt{BindingDB\_All\_202608} release.  Filters are applied in a fixed order
to a shrinking row set, so the removals and the retained count sum exactly to
the input total.
%%REV-END
We retained only finite, exact numerical $K_d$ measurements in nM for valid,
single-chain proteins of at most 1{,}200 residues and RDKit-valid ligands.
Inequality-qualified values were excluded.
%%REV-START
The dominant reduction is a measurement-type filter rather than a novelty
filter: the great majority of BindingDB rows report $K_i$, $\mathrm{IC}_{50}$ or
$\mathrm{EC}_{50}$ and carry no $K_d$ value at all, and are removed by the first
step.  The Davis-overlap exclusions, by contrast, remove only a small number of
rows.
%%REV-END
Before replicate reconciliation, any
external pair was excluded when its canonical isomeric
SMILES occurred anywhere in Davis or when its full normalised amino-acid
sequence occurred anywhere in Davis.  The resulting cohort is therefore
double-novel by exact ligand and target-sequence identity, although it is not a
claim of scaffold or remote-homology novelty. Measurements were then grouped by
canonical ligand--sequence pair. A pair with replicate measurements on opposite
sides of the \SI{30}{nM} boundary was removed (a pair-level conflict); for every
remaining pair, the median exact $K_d$ was retained and binarised as
$K_d<\SI{30}{nM}$. The complete row-to-pair accounting is reported in
Table~S13.

The released \texttt{best\_random} checkpoint was evaluated without retraining,
fine-tuning, calibration, or threshold selection on BindingDB.  We used the
existing random-validation threshold of 0.5959881544 and the same
640-dimensional ESM-2 representation used for Davis. Only newly extracted
external protein embeddings that passed dimension and sequence-length checks
were supplied to the frozen model.
AUROC and AUPRC confidence intervals use 600 pair-stratified bootstrap
replicates with seed 42.  Curation and frozen-inference provenance, including
source and prediction hashes, are reported in Table~S13.

%%REV-START
\subsection{Applicability-Domain Analysis}
\label{sec:ad_methods}

To interpret the external result, each curated BindingDB entity is placed
relative to the Davis training distribution on two axes.  On the ligand axis we
compute the maximum Tanimoto similarity between the non-chiral radius-2,
1024-bit Morgan fingerprint of the external compound and those of all 68 Davis
compounds.  On the target axis we align each external sequence against all 442
Davis sequences with the Biopython \texttt{Align.PairwiseAligner} in global mode
(match $=1$, mismatch $=0$, gap opening $=-1$, gap extension $=-0.5$) and retain
the maximum identity, defined as exact-residue matches divided by full alignment
length.  This is the same identity definition used for the within-Davis
cold-target audit, so the internal and external analyses are directly
comparable.  External pairs are then stratified by each axis, and AUROC, AUPRC,
positive counts and mean predicted probability are reported per stratum.  A pair
is treated as jointly inside the Davis domain when both scores are at least
0.30.  Calibration is assessed by comparing the frozen-threshold exceedance rate
and the predicted-probability distribution with the observed positive
prevalence.  The same two nearest-neighbour scores are computed live by the
interactive web server for each user-submitted pair.
%%REV-END

\subsection{Data Splitting Protocols and Similarity Audit}

\begin{itemize}[leftmargin=*]
\item \textbf{Random split} (70/10/20): drug--target pairs are assigned randomly;
  shared entities may appear across splits.
\item \textbf{Drug-ID-held-out split}: drugs are partitioned such that all interactions
  of a given compound identifier appear in exactly one split.  This prevents
  drug-identifier overlap but does not impose a scaffold-similarity threshold.
\item \textbf{Target-ID-held-out split}: targets are partitioned analogously; test
  target identifiers are absent from training.  This initial entity-held-out
  protocol does not impose a sequence-similarity threshold.
\end{itemize}

For the Drug-ID-held-out audit, we computed non-chiral radius-2, 1024-bit Morgan
fingerprints with RDKit and
reported the maximum Tanimoto similarity between each test compound and the
training compounds. For the Target-ID-held-out audit, we used RapidFuzz
\texttt{fuzz.ratio}, divided by 100, where normalised Indel similarity is
$1-d_\mathrm{Indel}/(|x|+|y|)$ and substitutions are represented by one
deletion plus one insertion. We computed the maximum value between each test
sequence and all training sequences and separately identified exact duplicate
sequences. The original
Target-ID-held-out partition contained six exact-sequence groups crossing the
training and test sets. We additionally constructed exact-sequence-grouped cold-target
and exact-sequence-grouped cold-both partitions, assigning all target identifiers with an identical
amino-acid sequence to the same split.
%%REV-START
These protocols use the same Davis records, 70/10/20 partition ratios, and
seed 42.  Grouping identical sequences removes exact duplicates only; it does
not cluster homologues, and Section~\ref{sec:discussion} states explicitly that
the resulting partitions are not a homology-reduced benchmark.
%%REV-END

For the strict cold-both protocol, the 70/10/20 ratios are applied independently
to drug identifiers and exact target-sequence groups. Only pairs for which both
entities belong to the same train, validation, or test block are retained;
cross-block combinations are excluded to prevent either entity from crossing a
partition. Consequently, the retained pair counts are 15{,}190/264/1{,}144,
and the pair-level proportions are not 70/10/20. Of the 30{,}056 Davis pairs,
13{,}458 cross-block pairs are not used in this cold-both evaluation.

To quantify residual target relatedness after exact-sequence grouping, we used
Biopython \texttt{Align.PairwiseAligner} in global mode (match $=1$, mismatch
$=0$, gap opening $=-1$, gap extension $=-0.5$) to align each
exact-sequence-grouped cold-target
test sequence against every training sequence. Identity is the number of
exact-residue matches divided by full alignment length; the maximum is retained
for each test target. We also implemented leakage-aware diagnostic controls on
this same exact-sequence-grouped cold-target partition. A drug-only control uses non-chiral
radius-2, 1024-bit
Morgan fingerprints in a logistic-regression classifier; a drug-promiscuity
control uses each drug's positive prevalence estimated only from training
targets; and a label-shuffled drug-only control permutes training labels while
leaving test labels unchanged.  Because every target identifier in this split
is unseen, a target-specific degree computed from its full label column would
leak test labels; the only target-only control that avoids that leakage assigns
the training positive prevalence uniformly.  These controls are diagnostic
rather than capacity-matched replacements for BioInteract.

\subsection{Evaluation Metrics}

For binary classification under severe class imbalance, we report AUROC and
AUPRC as primary metrics---AUPRC is particularly informative because it measures
positive-class precision across all recall levels. For every reported protocol,
F1, precision, and recall use a decision threshold selected once on the
corresponding validation precision--recall curve; AUROC and AUPRC remain
threshold-free.

\subsection{Implementation}

BioInteract is implemented in Python~3.10.14 using PyTorch~2.4.0+cu121 and
PyTorch Geometric~2.6.1.\cite{fey2019fast} Protein embeddings are extracted using
ESM-2\cite{lin2023evolutionary} (\texttt{esm2\_t30\_150M\_UR50D}, 640
dimensions; \texttt{fair-esm}~2.0.0). Molecular featurisation is performed with
RDKit~2025.03.6.\cite{rdkit} Revision diagnostics additionally use
Biopython~1.87 and RapidFuzz~3.14.5. Training employs automatic mixed precision (AMP) for
memory efficiency on a single NVIDIA GPU. Final checkpoint evaluation is run
under \texttt{torch.inference\_mode()} in full precision, using each checkpoint's
embedded architecture configuration.

%%====================================================================
\section{Results}
\label{sec:results}
%%====================================================================

%%REV-START
\subsection{Matched-Protocol Comparison with Baseline Methods}
\label{sec:matched}

Table~\ref{tab:comparison} and Fig.~\ref{fig:comparison} report the
three-model ladder of Section~\ref{sec:matched_methods}, each row trained by the
same harness on the same partitions.  All values in the table were produced in
this study; none is quoted from another publication.  Supporting Information
Table~S14 places the same ladder alongside two further architecture families
evaluated by the same harness on the random split.

\begin{table}[htbp]
\centering
\caption{Matched-protocol comparison on the Davis dataset (binary
  classification, $K_d < \SI{30}{nM}$).  Every model was trained in this study
  under one shared partition, seed, optimiser, schedule, loss and
  validation-selected threshold; no value is quoted from another publication.
  ``PLM'' indicates whether the model consumes the frozen
  \texttt{esm2\_t30\_150M\_UR50D} residue representation used by BioInteract.
  Best value in each column is in \textbf{bold}.  The Drug-ID-held-out columns
  are shown for completeness but are not used to rank the models: the shared
  early-stopping patience halted every run on that protocol at a first-epoch
  checkpoint (Section~\ref{sec:matched}).  Per-run provenance is given in
  Supporting Information Table~S16.}
\label{tab:comparison}
\resizebox{\textwidth}{!}{%
\begin{tabular}{@{}llccccccc@{}}
\toprule
 & & & \multicolumn{2}{c}{\textbf{Random}} &
     \multicolumn{2}{c}{\textbf{Target-ID-held-out}} &
     \multicolumn{2}{c}{\textbf{Drug-ID-held-out}} \\
\cmidrule(lr){4-5}\cmidrule(lr){6-7}\cmidrule(lr){8-9}
\textbf{Model} & \textbf{Interaction paradigm} & \textbf{PLM} & AUROC & AUPRC & AUROC & AUPRC & AUROC & AUPRC \\
\midrule
%%MATCHED-ROWS%%
\bottomrule
\end{tabular}
}
\end{table}

\begin{figure}[htbp]
\centering
\includegraphics[width=\linewidth]{figures/fig_comparison.pdf}
\caption{\textbf{Matched-protocol baseline comparison.} Panel A shows test
  AUROC and panel B shows test AUPRC for each row of the ladder under the random,
  Target-ID-held-out and Drug-ID-held-out protocols. All values were measured in
  this study under one shared training protocol. The Drug-ID-held-out bars are
  shown for completeness only and are not used to rank the models
  (Section~\ref{sec:matched}).}
\label{fig:comparison}
\end{figure}

\FloatBarrier

Both steps of the ladder improve on the step below them, but by amounts that
depend strongly on the splitting protocol.

One column should be set aside before reading the table.  Under the
Drug-ID-held-out protocol, validation AUROC dips below its first-epoch value for
several epochs in every run, so the shared early-stopping patience of six epochs
halted all three models before they recovered and restored a first-epoch
checkpoint.  That column therefore reflects the harness setting rather than the
converged behaviour of the three architectures, and it is not used to rank them;
the archived Drug-ID-held-out checkpoint reported in Table~S5, trained with a
longer patience, reaches AUROC 0.733.  The two remaining columns are read below.

Replacing the convolutional protein encoder with the frozen ESM-2 representation,
holding the drug encoder and the training recipe fixed, changes AUROC by
$+0.022$ and $+0.022$ and AUPRC by $+0.120$ and $+0.063$ under the random and
Target-ID-held-out protocols.  The contribution of the protein language model is
therefore consistent in direction across both usable protocols, and larger on
AUPRC than on AUROC.

Replacing concatenation with bidirectional cross-attention, holding both
encoders fixed, changes AUROC by $+0.006$ and $+0.056$ and AUPRC by $+0.067$ and
$+0.201$ under the same two protocols.  The interaction module is thus the
smaller of the two effects on the random split and the larger one when target
identifiers are held out, where it accounts for most of BioInteract's margin.
We read this as the clearest positive result of the matched comparison, subject
to the caveat that each configuration was trained once, so differences of a few
thousandths are not interpreted.

The Target-ID-held-out column is therefore where the architecture reported here
separates most clearly from its own controls: BioInteract reaches AUROC 0.934
and AUPRC 0.569 against 0.878 and 0.369 for the same encoders without
cross-attention, and 0.856 and 0.306 without the protein language model as well.
Because this is the protocol under which a model would be asked to score a newly
characterised target, we regard it, rather than the random split, as the
informative column of Table~\ref{tab:comparison}.  We do not set these values
against numbers published for other models: such a comparison requires a shared
partitioning, binarisation, threshold rule and model-selection rule, and no
controlled mapping between the two exists.

The random split discriminates poorly between architectures, and the breadth
check in Supporting Information Table~S14 shows why this matters.  Under the
same harness on that split, the five architectures tested span 0.885 to 0.913
AUROC, a range narrow enough that the ordering within it carries little
information.  The two additional families were not run on the entity-held-out
protocols, so we make no claim about how they behave there.
What the random-split breadth check does establish is that the protocol on which
DTI models are most often compared is the one on which architectural differences
are least visible, and that a favourable random-split number should not be read
as evidence for a particular architecture.  The results reported here are
consistent with that view rather than in tension with it: BioInteract's
advantage within the controlled ladder appears where the evaluation is harder,
not where it is easiest.
%%REV-END

\subsection{Ablation Study}
\label{sec:ablation}

We assessed the contribution of the principal architectural components by
removing one component at a time from the full model (Fig.~\ref{fig:ablation}).
%%REV-START
These ablation variants come from the archived experiment series rather than
from the matched-protocol harness, and their absolute values are therefore not
interchangeable with Table~\ref{tab:comparison}.  They are retained because
they cover two components---the domain-label channel and graph
augmentation---that the matched-protocol contrasts do not isolate.
%%REV-END

\begin{figure}[htbp]
\centering
\includegraphics[width=\linewidth]{figures/fig_ablation.pdf}
\caption{\textbf{Ablation-study results.} Panel A shows the random- and
  Target-ID-held-out metric matrix for the full model and for each variant with
  one component removed. Panel B reports the AUROC change after removal
  of cross-attention, the domain-label channel, ESM-2, or graph augmentation.
  ESM-2 removal produces the largest AUROC decrease, followed by removal of
  cross-attention. Values are from the archived experiment series and are not
  on the same footing as the matched-protocol results in
  Table~\ref{tab:comparison}.}
\label{fig:ablation}
\end{figure}

\FloatBarrier

Removing ESM-2 changes AUROC by $-0.059$ and $-0.137$ for the random and
Target-ID-held-out splits, respectively; removing cross-attention changes AUROC
by $-0.040$ and $-0.079$. The domain-label and graph-augmentation variants show
small mixed changes ($+0.002/-0.018$ and $+0.004/-0.004$, respectively). These
ablation results describe the tested Davis configuration and do not establish
structural binding mechanisms or out-of-family generalisation.

\subsection{Strict Within-Davis Cold-Start Evaluation}

The maximum normalised Indel similarity between each Target-ID-held-out test
sequence and the original
training targets ranges from 0.382 to 1.000 (median 0.579). The audit also found
18 duplicate-sequence groups involving 81 target
identifiers in the full dataset and six groups crossing the original train/test
boundary.  We therefore report the exact-sequence-grouped results separately
and limit conclusions to the observed Davis benchmark performance.

\begin{table}[htbp]
\centering
\caption{Strict within-Davis cold-start evaluation after grouping target
  identifiers with identical amino-acid sequences before partitioning.  Both
  protocols use the original records, seed 42, and the same preprocessing and
  architecture, but are separately trained from initialisation. For cold-both,
  70/10/20 is applied independently to drug IDs and
  target-sequence groups and only within-block pairs are retained. AUROC and
  AUPRC are threshold-free test metrics.}
\label{tab:strict_results}
\begin{tabular}{l@{\hspace{1em}}rrrr}
\toprule
\textbf{Strict protocol} & \textbf{Test pairs} & \textbf{Test positives} & \textbf{AUROC} & \textbf{AUPRC} \\
\midrule
Exact-sequence-grouped cold-target & 5{,}984 & 208 & 0.873 & 0.383 \\
Exact-sequence-grouped cold-both   & 1{,}144 & 38  & 0.552 & 0.038 \\
\bottomrule
\end{tabular}
\end{table}

After grouping identical target sequences, the exact-sequence-grouped cold-target test gives
AUROC~$=0.873$ and AUPRC~$=0.383$ (Table~\ref{tab:strict_results}).  Requiring
both unseen drug identifiers and exact-sequence-grouped targets is substantially
harder (AUROC~$=0.552$, AUPRC~$=0.038$); this test contains 38 positive pairs.
These values are reported as an informative within-Davis stress test, not as a
claim of out-of-family or structure-level generalisation.  Threshold-dependent
metrics for these small positive sets are secondary and use a threshold selected
only on the corresponding validation set.

The cold-both train, validation, and test blocks retain 15{,}190, 264, and
1{,}144 pairs, respectively; 13{,}458 cross-block pairs are excluded by the
double-held-out construction. The validation block contains only 10 positives.
At its validation-selected threshold (0.936), the cold-both test yields F1,
precision, and recall of 0.000, underscoring that its near-random AUROC and low
AUPRC do not support a positive generalisation claim. Full counts and
threshold-dependent metrics are reported in Table~S7.

\begin{table}[htbp]
\centering
\caption{Leakage-aware diagnostic controls on the exact-sequence-grouped
  cold-target test. All controls use only training-partition labels. They
  quantify available drug-side signal; they are not capacity-matched substitutes
  for BioInteract.}
\label{tab:diagnostic_controls}
\begin{tabular}{lrr}
\toprule
\textbf{Model or control} & \textbf{AUROC} & \textbf{AUPRC} \\
\midrule
BioInteract (exact-sequence-grouped cold-target) & 0.873 & 0.383 \\
Drug-only Morgan logistic         & 0.764 & 0.164 \\
Drug-promiscuity prior            & 0.764 & 0.163 \\
Target-only constant prior        & 0.500 & 0.035 \\
Label-shuffled drug-only          & 0.407 & 0.029 \\
\bottomrule
\end{tabular}
\end{table}

The drug-only and drug-promiscuity controls reach AUROC values near 0.764,
showing that ligand-side activity or promiscuity signal contributes materially
to this within-kinome task (Table~\ref{tab:diagnostic_controls}).  BioInteract
exceeds these diagnostic controls, but the comparison does not prove that its
residue attributions recover physical interactions.  The label-shuffled control
is near or below the positive-prevalence baseline, confirming that the drug-only
signal depends on the observed training labels rather than fingerprint ordering.

%%REV-START
For the exact-sequence-grouped cold-target test, the maximum global sequence
identity to any training target ranges from 0.231 to 0.867 (median 0.494).
Figure~\ref{fig:identity_bins} stratifies the 88 held-out targets by that
maximum identity and reports, for each stratum, BioInteract's AUROC and AUPRC
together with the number of test targets, test pairs and positive pairs that
support them.  Performance is not monotone in identity: the lowest-identity
stratum gives the lowest AUROC (0.806), the two middle strata are the strongest
(0.908 in both), and the highest-identity stratum falls back to 0.862 with an
AUPRC of 0.121 computed from only three positive pairs.  The per-stratum
positive counts make clear that the two outer strata are the least stable, and
the analysis therefore documents the residual relatedness distribution rather
than demonstrating robust transfer to homologously distant targets.  Numerical
values are tabulated in Supporting Information Table~S15.
%%REV-END

\begin{figure}[htbp]
\centering
\includegraphics[width=0.86\linewidth]{figures/fig_identity_bins.pdf}
\caption{\textbf{Identity-stratified performance on the exact-sequence-grouped
  cold-target test.} Held-out targets are grouped by their maximum full-length
  global sequence identity to any training target. Bars give BioInteract's
  AUROC and AUPRC within each stratum; the annotation under each stratum gives
  the number of held-out targets, test pairs and positive pairs contributing to
  it. The dashed line marks the positive prevalence of the full cold-target
  test. Strata with few positives yield unstable AUPRC estimates.}
\label{fig:identity_bins}
\end{figure}

\FloatBarrier

Under the Drug-ID-held-out split, the maximum test-to-training Morgan Tanimoto
similarity ranged from 0.196 to 0.538 (median 0.257).  The fingerprint audit is
included to clarify that Drug-ID-held-out means unseen compound identifiers, not
an externally validated scaffold-hopping or out-of-distribution guarantee.

\subsection{Strict External BindingDB Evaluation}

%%REV-START
\subsubsection{Reproducible Construction of the External Cohort}

Table~\ref{tab:bindingdb_curation} gives the complete row-level accounting for
the curation described in Section~\ref{sec:bindingdb_methods}.  The
curated-articles snapshot on which the pre-specified evaluation is based
contains 93{,}712 rows, of which 91{,}062 (\SI{97.2}{\percent}) report no $K_d$
at all because they record $K_i$, $\mathrm{IC}_{50}$ or $\mathrm{EC}_{50}$
instead.  Only 2{,}650 rows carry any $K_d$ value.  Every subsequent filter,
including both Davis-overlap exclusions, removes at most a few hundred further
rows, so the size of the final cohort is determined almost entirely by the
requirement of an exact $K_d$ measurement in a snapshot of that size, not by
the novelty criteria.

The same pipeline applied to the complete \texttt{BindingDB\_All\_202608}
release makes this explicit: from 3{,}234{,}499 rows the identical filters
retain 52{,}085 eligible measurements covering 27{,}079 ligands and 3{,}472
distinct target sequences, a cohort roughly 24 times larger than the
curated-articles cohort.  The filter definitions therefore scale as expected,
and the modest size of the pre-specified cohort is a property of the snapshot
that was fixed before the evaluation rather than of the criteria themselves.
Section~\ref{sec:bindingdb_full} reports the frozen evaluation on that larger
cohort.  Both waterfalls are reproduced by a released audit script, and the
per-snapshot archive hashes are given in Supporting Information Table~S17.

\begin{table}[htbp]
\centering
\small
\caption{Row-level filtering waterfall for the external BindingDB cohorts.
  Filters are applied in the listed order to a shrinking row set, so removals
  and the retained total sum exactly to the input rows.  Supporting Information
  Table~S13 reports the same curation for the curated-articles snapshot with the
  single-chain test applied first; the ordering changes how rows are attributed
  between the first three filters but not the eligible total.}
\label{tab:bindingdb_curation}
\begin{tabular}{lrrrr}
\toprule
 & \multicolumn{2}{c}{\textbf{Curated articles, 202607}} &
   \multicolumn{2}{c}{\textbf{Complete release, 202608}} \\
\cmidrule(lr){2-3}\cmidrule(lr){4-5}
\textbf{Filter} & Removed & Remaining & Removed & Remaining \\
\midrule
Input rows                                 & ---     & 93{,}712 & ---           & 3{,}234{,}499 \\
No $K_d$ reported ($K_i$/IC$_{50}$/EC$_{50}$ rows) & 91{,}062 & 2{,}650 & 3{,}104{,}331 & 130{,}168 \\
$K_d$ inequality-qualified or unparsable   & 160     & 2{,}490  & 47{,}860      & 82{,}308 \\
Target not a single protein chain          & 27      & 2{,}463  & 2{,}157       & 80{,}151 \\
Ligand SMILES not parsable by RDKit        & 8       & 2{,}455  & 93            & 80{,}058 \\
Non-standard residues or sequence $>$ 1{,}200 aa & 141 & 2{,}314 & 7{,}644      & 72{,}414 \\
Ligand identical to a Davis compound       & 62      & 2{,}252  & 5{,}738       & 66{,}676 \\
Sequence identical to a Davis target       & 94      & 2{,}158  & 14{,}591      & 52{,}085 \\
\midrule
Eligible measurements                      & \multicolumn{2}{c}{2{,}158} & \multicolumn{2}{c}{52{,}085} \\
Unique eligible ligands / sequences        & \multicolumn{2}{c}{1{,}197 / 400} & \multicolumn{2}{c}{27{,}079 / 3{,}472} \\
\bottomrule
\end{tabular}
\end{table}
%%REV-END

We next performed the pre-specified, frozen external evaluation on the curated
BindingDB cohort.  After all measurement-quality, exact-identity novelty, and
replicate-consistency filters, it contains 1{,}932 pairs from 1{,}195 ligands
and 396 targets: 516 high-affinity and 1{,}416 low-affinity pairs
(Table~\ref{tab:bindingdb_external}).  The model was neither retrained nor
recalibrated, and no BindingDB label was used to choose the checkpoint or
threshold.

The frozen checkpoint achieved AUROC $=0.560$ (95\% CI 0.531--0.589) and
AUPRC $=0.340$ (95\% CI 0.314--0.372).  At the frozen Davis threshold, F1 was
0.070, precision was 0.731, and recall was 0.037.  Thus, although the exact
double-novel cohort is a valid external test, it does not support a broad DTI-transfer claim.
Instead, it shows limited transfer from this
Davis-trained classifier to the broader BindingDB cohort.  This external result
does not alter the within-Davis checkpoint reconstructions and is not used to
interpret attention, residue attribution, binding pockets, or mechanisms.

%%REV-START
\subsection{Frozen Evaluation on the Complete BindingDB Release}
\label{sec:bindingdb_full}

Because the pre-specified cohort is small, we repeated the frozen evaluation on
the much larger cohort obtained by applying the identical protocol to the
complete \texttt{BindingDB\_All\_202608} release
(Table~\ref{tab:bindingdb_curation}).  After replicate reconciliation this
cohort contains 38{,}193 double-novel pairs from 26{,}433 ligands and 3{,}462
target sequences, of which 13{,}264 are high-affinity.  The same
\texttt{best\_random} checkpoint, the same frozen threshold, and the same
freshly extracted \texttt{fair-esm} representation were used; no BindingDB label
influenced the checkpoint, the threshold, or any preprocessing decision.  This
evaluation was specified and run after the reviewers' comments, and is reported
as a verification of the pre-specified result rather than as a replacement
for it.

%%FULLDB-RESULT%%

\subsection{Applicability Domain of the External Cohort}
\label{sec:applicability}

The external cohort was constructed to exclude exact Davis entities, so it is
novel by construction; the relevant question is how far novel it is, and whether
distance from the Davis training distribution accounts for the loss of
performance.  Table~\ref{tab:applicability} places every external pair on the
two axes defined in Section~\ref{sec:ad_methods}.

The cohort is overwhelmingly remote from Davis on both axes.  The median
external ligand has a maximum Tanimoto similarity of 0.197 to any of the 68
Davis compounds, and \SI{92.6}{\percent} of external ligands fall below 0.30.
The median external target has a maximum global sequence identity of 0.238 to
any of the 442 Davis kinases, and \SI{90.2}{\percent} fall below 0.30.  This is
expected---Davis contains a single protein family---but it means the headline
external AUROC is dominated by pairs far outside the training distribution.

Performance tracks that distance closely.  In the two lowest ligand-similarity
strata, which contain \SI{92.8}{\percent} of the pairs, AUROC is 0.509 and
0.504, that is, indistinguishable from chance.  In the two strata at or above
0.30 it rises to 0.700 and 0.723.  The target axis behaves the same way: AUROC
is 0.538 and 0.517 below 0.30 identity and 0.749 and 0.623 above it.  Restricted
to the 117 pairs that are simultaneously at or above 0.30 on both axes, the
frozen checkpoint gives AUROC~$=0.596$ and AUPRC~$=0.719$ against a positive
prevalence of 0.641 in that subset.  The in-domain signal is therefore real but
modest, and the out-of-domain signal is absent.

\begin{table}[htbp]
\centering
\small
\caption{Applicability-domain stratification of the strict external BindingDB
  cohort. Each pair is placed by the maximum Tanimoto similarity of its ligand
  to the 68 Davis compounds and by the maximum global sequence identity of its
  target to the 442 Davis targets. ``Entities'' counts the distinct external
  ligands or targets in the stratum. Mean $\hat{p}$ is the mean predicted
  probability from the frozen checkpoint.}
\label{tab:applicability}
\begin{tabular}{lrrrrrr}
\toprule
\textbf{Stratum} & \textbf{Pairs} & \textbf{Positives} & \textbf{Entities} &
\textbf{AUROC} & \textbf{AUPRC} & \textbf{Mean $\hat{p}$} \\
\midrule
\multicolumn{7}{l}{\emph{Maximum ligand Tanimoto similarity to Davis compounds}} \\
$<0.20$          & 1{,}076 & 296 & 627 & 0.509 & 0.305 & 0.375 \\
$0.20$--$<0.30$  &    717 & 144 & 479 & 0.504 & 0.217 & 0.366 \\
$0.30$--$<0.40$  &     46 &  19 &  36 & 0.700 & 0.614 & 0.420 \\
$\geq 0.40$      &     93 &  57 &  53 & 0.723 & 0.785 & 0.444 \\
\midrule
\multicolumn{7}{l}{\emph{Maximum global sequence identity to Davis targets}} \\
$<0.20$          &     14 &   1 &   7 & 0.538 & 0.143 & 0.397 \\
$0.20$--$<0.30$  & 1{,}721 & 415 & 350 & 0.517 & 0.266 & 0.371 \\
$0.30$--$<0.40$  &     73 &  37 &  15 & 0.749 & 0.722 & 0.409 \\
$\geq 0.40$      &    124 &  63 &  24 & 0.623 & 0.647 & 0.429 \\
\midrule
Both axes $\geq 0.30$ & 117 & 75 & --- & 0.596 & 0.719 & --- \\
\bottomrule
\end{tabular}
\end{table}

A second, independent failure mode is visible in the score distribution rather
than in the ranking.  The frozen Davis threshold of 0.596 is exceeded by only
\SI{1.3}{\percent} of external pairs, although \SI{26.7}{\percent} of them are
positive, which is why the frozen-threshold recall is 0.037 while precision
remains 0.731.  The predicted probabilities are compressed into a narrow
band---median 0.363, 95th percentile 0.437---with a mean of 0.392 for positives
against 0.370 for negatives.  The operating point transferred from Davis is
therefore miscalibrated for this cohort even where the ranking retains
information.

Assay and label heterogeneity, by contrast, does not appear to be a major
contributor.  Only 143 of the 1{,}932 retained pairs had replicate
measurements, and only 15 candidate pairs were discarded because replicate
measurements straddled the \SI{30}{nM} boundary, a pair-level conflict rate of
\SI{0.8}{\percent}.

The same two nearest-neighbour scores are now computed live by the interactive
web server for every submitted ligand--target pair and reported alongside the
classifier score with an explicit inside/borderline/outside assessment
(Fig.~\ref{fig:webserver}), so that a user is told when a query lies in the
regime where the external evaluation showed chance-level ranking.
%%REV-END

\subsection{Training Dynamics}

Fig.~\ref{fig:training} reports a single checkpoint-matched, contiguous
validation-AUROC and training-loss trace for each archived protocol. The
duplicated Random trace and two restarted Target-ID fragments were excluded
because their epoch counters reset or their peak validation AUROC did not match
the released checkpoint. These plots describe observed optimisation paths only;
they do not establish absence of overfitting, quantify validation-loss behaviour,
or estimate repeated-seed variability.

\begin{figure}[htbp]
\centering
\includegraphics[width=\textwidth]{figures/fig4_training.pdf}
\caption{\textbf{Training dynamics across splitting protocols.}
  Panel A: validation AUROC over training epochs.  Panel B: training loss curves.
  Each line is one contiguous archived trace whose peak validation AUROC agrees
  with the saved checkpoint; duplicated and restarted segments are not plotted.
  No fitted or interpolated trends are shown. The trajectories do not establish
  absence of overfitting, validation-loss behaviour, or repeated-seed variability.}
\label{fig:training}
\end{figure}

\FloatBarrier

\subsection{Interpretability: Cross-Attention Analysis}
\label{sec:interpretability}

All interpretability figures and tables are derived directly from the saved
case-level outputs in the project interpretability-results directory, including
the report JSON and the corresponding heatmap, profile, and Grad-CAM renderings.
Accordingly, we report the model's actual residue-level attention and atom-level
Grad-CAM outputs rather than schematic placeholders.  These are attribution
outputs of the trained model, not independently validated residue contacts or
functional-group assignments.

\subsubsection{Global Attention Sparsity}

Analysis of cross-attention patterns across all 1{,}506 positive Davis
drug--target pairs
reveals a global attention sparsity of \SI{99.2}{\percent}
(Table~\ref{tab:global_attn}, Fig.~\ref{fig:sparsity}) under the stated
normalisation threshold. Figure~\ref{fig:sparsity} is computed directly from
the archived 1{,}301{,}380 pair-normalised residue scores, not from summary
bars; its top 1\% and 5\% of scores account for 78.9\% and 90.4\% of their
total attribution mass, respectively. This describes concentration of model
attribution; it does not demonstrate the size or location of a physical binding
interface.

\begin{table}[htbp]
\centering
\caption{Global cross-attention statistics computed across all 1{,}506 positive
  Davis pairs using \texttt{checkpoints/best.pt}. The values quantify the
  concentration of the model's pair-normalised
  residue-attribution scores; they do not provide structural validation of a
  binding pocket.}
\label{tab:global_attn}
\begin{tabular}{p{0.72\linewidth}c}
\toprule
\textbf{Metric} & \textbf{Value} \\
\midrule
Mean residue attention      & 0.0044 \\
Median residue attention    & 0.0003 \\
Top-1\% threshold           & 0.0572 \\
Top-5\% threshold           & 0.0036 \\
Fraction of residues with normalised attribution $<0.1$ of the pair-specific maximum & 99.2\% \\
Top 1\% / top 5\% cumulative attribution mass & 78.9\% / 90.4\% \\
Samples analysed            & 1{,}506 \\
Residue-score observations   & 1{,}301{,}380 \\
\bottomrule
\end{tabular}
\end{table}

\begin{figure}[htbp]
\centering
\includegraphics[width=\linewidth]{figures/fig5_attention_sparsity.pdf}
\caption{\textbf{Attention sparsity analysis.}
  (A)~Histogram of the archived pair-normalised per-residue attention scores
  across all 1{,}506 positive Davis pairs (1{,}301{,}380 observations; logarithmic
  axes); the dashed line denotes the 0.1 sparsity threshold. (B)~Cumulative
  attribution mass after sorting those same scores in descending order. The
  \SI{99.2}{\percent} value is an attribution-concentration statistic, not a
  claim about a contiguous protein-surface region.}
\label{fig:sparsity}
\end{figure}

\FloatBarrier

Although softmax attention is dense before normalisation, the observed
concentration arises after training on the limited and imbalanced Davis data.
It may reflect learned regularities, data composition, or both; it is therefore
not presented as independent evidence that the model has recovered molecular
recognition principles.

\subsubsection{Input-Integrity Audit of Nominal ABL1 Variant Records}
\label{sec:abl1}

%%REV-START
This audit tests one falsifiable statement about the benchmark inputs rather
than about the model: \emph{if} the Davis records labelled with ABL1
point-mutation names encoded the corresponding sequence substitutions, \emph{then}
the amino-acid sequences supplied to the model for those records would differ
from the wild-type sequence and from one another.  The evaluation criterion is
therefore exact: for every nominal ABL1 record we compare the full amino-acid
string and its SHA-256 digest in \texttt{target\_sequences.csv}.

The audit falsifies that statement.  Fifteen target records whose nominal names
begin with ABL1---including the wild-type and the named F317I, F317L, M351T and
E255K records---all carry the same 1{,}167-residue amino-acid sequence and the
same SHA-256 digest (Supporting Information, Table~S9).  The consequence is a
property of the benchmark, not a finding about mutation biology: because
variant-specific sequence information is absent from the model inputs, any
difference in predicted score between two nominal ABL1 variants can only reflect
the paired ligand, and no experiment on these records can measure
mutation-specific or resistance-specific model behaviour.  The value of the
audit is that it defines a boundary on what the Davis benchmark can be used to
test, and it is reported for that reason alone; the mutant-comparison analysis
that appeared in the original submission is not part of this manuscript.
%%REV-END

\subsection{Case Study: Functional-Group Attribution}

The saved Grad-CAM outputs for the two archived non-mutant Davis cases provide a
compact example of atom-level scores aggregated to functional groups
(Table~\ref{tab:functional_group}).  The entries are model-native attribution
scores rather than measured contacts or interaction types.

%%REV-START
\begin{table}[htbp]
\centering
\caption{Model-native functional-group attribution for the two archived
  non-mutant Davis cases. Scores are normalised atom-level Grad-CAM values
  aggregated by functional group; they do not identify physical contacts. A
  dash marks a group with no RDKit substructure match in that ligand.}
\label{tab:functional_group}
\begin{tabular}{lcc}
\toprule
\textbf{Functional group} & \textbf{EGFR--D0014} & \textbf{BRAF--D0023} \\
 & (class score 0.995) & (class score 0.099) \\
\midrule
Carbonyl      & 0.431 & --- \\
Amide         & 0.389 & --- \\
Ether         & 0.136 & --- \\
Halogen       & 0.134 & --- \\
Amino         & 0.090 & --- \\
Aromatic ring & 0.000 & 0.059 \\
Heterocycle N & 0.000 & 0.000 \\
Hydroxyl      & ---   & 0.000 \\
\bottomrule
\end{tabular}
\end{table}
%%REV-END

\FloatBarrier

The two cases differ markedly in classifier score and attribution profile.
EGFR--D0014 has a high-affinity-class score of 0.995 and its largest aggregated
groups are carbonyl (0.431) and amide (0.389).  BRAF--D0023 has a score of 0.099
and only small saved values for aromatic ring (0.059) and hydroxyl (0.000),
despite its archived positive affinity of 0.19 nM. We therefore treat
BRAF--D0023 as a low-scoring failure-mode example for an experimentally positive
pair. This is a descriptive limitation of model behaviour, not an explanation of
target binding or an assessment of functional-group chemistry.

%%====================================================================
\section{Discussion}
\label{sec:discussion}
%%====================================================================

\subsection{From Prediction Accuracy to Inspectable Hypotheses}

BioInteract is intended to make model behaviour more inspectable than a purely
scalar DTI prediction.  Its attention matrices and Grad-CAM scores can generate
hypotheses about which representations contributed to a prediction.  Such
hypotheses can subsequently be assessed with crystal structures, mutagenesis,
or structure--activity relationship (SAR) experiments; the present study does
not replace these measurements.

%%REV-START
The matched-protocol comparison sharpens what can and cannot be concluded from
the Davis benchmark, and the two conclusions point in opposite directions.

The favourable one is that, within a ladder where a single component changes at
a time, both components earn their place, and the cross-modal interaction module
earns it where the evaluation is hardest.  Cross-attention adds 0.201 AUPRC over
concatenation under the Target-ID-held-out protocol against 0.067 under the
random split, with both encoders and the entire training recipe held fixed.  A
comparison run only on the random split would have judged that module nearly
worthless.

The qualifying one is that the random split, on which DTI architectures are most
often ranked, barely separates them: five architectures trained under the same
harness span 0.885 to 0.913 AUROC (Supporting Information Table~S14).  We
therefore treat random-split ordering as weak evidence about architecture, and
we do not rest any conclusion on it.

We therefore present BioInteract as one audited point in this design space
rather than as a state-of-the-art claim.  These Davis results
do not show that BioInteract maps a physical binding pocket, assigns an
interaction type to an atom pair, ranks a large compound library, or assesses
treatment resistance.
%%REV-END

\subsection{The 99.2\% Attention Sparsity}

The observed \SI{99.2}{\percent} attention sparsity is a property of the
pair-normalised, residue-pooled attribution scores under the chosen threshold. Although it is
compatible with a model focusing on a small subset of residues, it can also be
affected by the kinase-focused and imbalanced training data.  It should not be
read as evidence that the highlighted residues constitute a true binding site.
%%REV-START
More generally, work that supervises non-covalent interaction labels directly
has shown that concentrated or plausible-looking attention is not by itself
evidence of recovered
interactions.\cite{li2020monn,gim2023arkdta,seo2024blendnet}  We accordingly
report the sparsity statistic as a description of how the trained model
distributes attribution mass, and not as support for an interpretability claim.
%%REV-END

\subsection{Scope of Cold-Start Evaluation}

%%REV-START
The exact-sequence-grouped cold-target and cold-both evaluations, which provide
a strict within-Davis stress test, involve kinase proteins and do not
demonstrate performance on unrelated protein families or pocket classes.
Equally important, exact-sequence grouping controls duplicate sequences only.
Held-out targets retain substantial residual relatedness to the training
set---the maximum global identity has a median of 0.494 and reaches 0.867---so
these partitions must not be read as a fully homology-reduced cold-target
benchmark.  A benchmark that clustered targets by kinase-domain similarity, or
that adopted a predefined clustered split at a fixed identity cut-off, would
provide the stronger control and is a useful direction for future evaluation.
%%REV-END

The Drug-ID-held-out result is likewise restricted to the chemical diversity of 68
kinase inhibitors.  The observed Morgan-similarity distribution documents the
nearest-neighbour relationship between held-out and training molecules, but it
does not establish scaffold hopping or out-of-distribution robustness.

\subsection{Diagnostic Controls and Remaining Dataset Bias}

The exact-sequence-grouped cold-target diagnostic controls demonstrate that a target-independent
drug representation retains substantial ranking signal in Davis.  The
drug-only Morgan and drug-promiscuity controls each achieve AUROC near 0.764,
whereas BioInteract reaches 0.873 under the same exact-sequence-grouped
partition.  This improvement is compatible with the use of target information,
but it neither excludes ligand-side activity bias nor establishes that the model
has learned a molecular-recognition mechanism.  We therefore interpret the
control comparison as a boundary on the Davis result, not as causal evidence for
any particular encoder or attribution map.  The label-shuffled control confirms
that the observed drug-only signal depends on the training labels; it is not a
test of structural interaction recovery.

%%REV-START
\subsection{Why External Transfer Is Limited}
\label{sec:transfer}

The external result in Section~\ref{sec:applicability} is best read as two
separate failures with different remedies.

The first is a domain failure.  More than nine-tenths of the external pairs sit
below a Tanimoto similarity of 0.30 to any Davis compound and below a sequence
identity of 0.30 to any Davis target, and it is exactly those pairs whose AUROC
is indistinguishable from chance.  Within the small in-domain subset, ranking
recovers to AUROC 0.60--0.75 depending on the axis.  A model trained on 68
inhibitors and one protein family cannot be expected to rank pairs drawn from
the whole of BindingDB, and the stratified result shows that the aggregate
AUROC of 0.560 is a mixture of a usable in-domain regime and an uninformative
out-of-domain one rather than a uniform level of weak performance.  The remedy
is broader training data, not a different attribution mechanism.

The second is a calibration failure that is independent of ranking.  The
decision threshold selected on the Davis validation set is exceeded by
\SI{1.3}{\percent} of external pairs even though \SI{26.7}{\percent} are
positive, and the predicted probabilities occupy a narrow band well below that
threshold.  Any deployment on a new distribution would therefore require the
threshold to be re-derived on labelled data from that distribution; reporting a
frozen-threshold F1 across distributions, as we do in
Table~\ref{tab:bindingdb_external}, understates ranking quality and overstates
the usefulness of the fixed operating point.

Two candidate explanations are not supported by the data.  Assay and label
heterogeneity is small in this cohort: the pair-level replicate-conflict rate is
\SI{0.8}{\percent}, and conflicting pairs were removed before evaluation.  And
the result is not an artefact of cohort size, since the same filters applied to
the complete BindingDB release retain 52{,}085 measurements over 3{,}472 target
sequences (Table~\ref{tab:bindingdb_curation}).  We therefore attribute the
limited transfer principally to chemical and target novelty relative to Davis,
compounded by threshold miscalibration, and we report the applicability-domain
scores in the released web server so that this boundary is visible at the point
of use.
%%REV-END

\subsection{Functional-Group Attribution Without Structural Supervision}

The Grad-CAM functional-group ranking is qualitatively compared with kinase
SAR literature, but it remains a gradient-based attribution.  In particular,
the analysis does not determine whether a group forms a hydrogen bond, a
pi-stacking interaction, a cation--pi interaction, or a van der Waals contact,
nor whether an implicated protein residue contributes through its backbone or
side chain.  These questions require a structural complex or targeted
experiments.\cite{zhang2009targeting}

\subsection{Future Directions}

The strict BindingDB evaluation directly indicates limited transfer beyond the
Davis benchmark; it therefore cannot support a protein-family-agnostic,
broad-pocket, or chemical-space generalisation claim.  Evaluation after
retraining and validation on a wider, pre-specified multi-family protocol is
required to establish whether the limitation arises from the Davis training
distribution, task definition, or model. Direct comparison with independent
experimental evidence is also needed before attributions can be assigned a
structural meaning.

%%REV-START
The most informative next step for the attribution outputs is not a further
qualitative case study but a supervised comparison.  Models trained or
regularised with explicit non-covalent interaction labels provide both the
benchmark design and the evaluation criterion
required.\cite{li2020monn,gim2023arkdta,seo2024blendnet}  A future structural
audit should use a non-overlapping benchmark of experimentally resolved
complexes, predefine the atom--residue contact criterion, and compare model
attributions with the resulting contact labels (for example, from
PLIP\cite{salentin2015plip}) before structural meaning is assigned.  No overlap
analysis between the Davis cases and PDBBind or LP-PDBBind, and no PLIP contact
validation, was performed in this study.
%%REV-END

The present implementation was audited on one NVIDIA GPU. Although the
2{,}442{,}083 trainable parameters do not increase with the number of training
pairs, data-processing time, graph batching, ESM-2 embedding extraction and
storage, and accelerator memory must be profiled at larger scale. Distributed
training was not evaluated, and no claim is made that the current implementation
scales linearly to BindingDB-sized data.

Evaluation on curated allosteric and cryptic-site benchmarks, with appropriate
structural or dynamic context, is likewise required before BioInteract can be
assessed for non-canonical drug--target interactions.  Existing work highlights
the need for explicit allosteric-site labels and benchmark design,\cite{wu2022allostericbenchmark,gatlin2026allostericgrammar}
and for structural or dynamic treatment of cryptic pockets.\cite{zhang2026crypticpockets}
CryptoBench provides a dedicated cryptic protein--ligand benchmark that can help
define such a future evaluation.\cite{skrhak2024cryptobench}

Extending BioInteract from binary interaction classification to
continuous binding-affinity regression ($pK_d$) would provide the finer-grained
potency estimates needed for lead optimisation, where accurate ranking of closely
related analogues is as important as binary activity prediction.

%%====================================================================
\section{Conclusion}
\label{sec:conclusion}
%%====================================================================

We have presented BioInteract, a DTI-prediction framework that combines
chemistry-aware molecular graph encoding, ESM-2/physicochemical residue
representations, and bidirectional cross-attention.  The model provides
prediction scores together with inspectable model-native attention and Grad-CAM
attributions.

%%REV-START
This revision states the evidence boundary of those outputs explicitly.  Under a
matched training protocol, both components of the three-model ladder contribute,
and the cross-modal interaction module contributes most under the
entity-held-out protocols rather than on the random split.  Because the random
split separates architectures only weakly, we make no claim of general
superiority and treat random-split rankings as weakly informative about
architecture.  The
observed attention concentration and functional-group ranking describe model
behaviour, and, in view of results obtained with explicit non-covalent
interaction supervision, we do not present them as recovered binding contacts or
mechanisms.  We distinguish Target-ID-held-out performance from
exact-sequence-grouped cold-start evaluation after detecting duplicate sequences
in the original target partition, while noting that neither is a
homology-reduced benchmark.  The strict external BindingDB test shows limited
transfer of the frozen Davis checkpoint, and the accompanying
applicability-domain analysis identifies the conditions under which the released
model should and should not be trusted.  Broader datasets, clustered splits and
direct structural validation remain necessary for generalisation and mechanistic
claims.
%%REV-END

BioInteract is available as a publicly accessible interactive web interface at
\url{https://huggingface.co/spaces/AI4deeperScience/BioInteract} for the
custom-prediction workflow shown in Fig.~\ref{fig:webserver}.  The repository
README documents the command-line workflow for the available attribution
outputs.

%%REV-START
\begin{figure}[htbp]
  \centering
  \includegraphics[width=0.995\linewidth]{figures/fig_webserver.pdf}
  \caption{\textbf{BioInteract interactive web server.}
    Panel A shows the input and result state of the publicly accessible
    interface at \url{https://huggingface.co/spaces/AI4deeperScience/BioInteract}
    for the released Davis example pair D0017 (dasatinib) and T0210 (LCK,
    509~residues), including the high-affinity-class classifier score and the
    applicability-domain readout.  For every submitted pair the interface now
    reports the ligand's maximum Tanimoto similarity and the target's maximum
    global sequence identity to the Davis training entities, together with an
    explicit inside/borderline/outside assessment.  Panels B and C show the two
    model-native attribution views returned for the same pair: the atom--residue
    attention map over the highest-attribution residue window, and the ranked
    top-ten residue attributions.  All three panels are rendered as vector
    graphics from the released pipeline.  These are model-native attribution
    views, not physical-contact or hotspot labels.  Command-line instructions
    for generating the available attribution outputs are provided in the public
    repository.}
  \label{fig:webserver}
\end{figure}
%%REV-END

BioInteract represents a step toward AI-assisted DTI prediction in which model
scores are accompanied by transparent attribution outputs and explicit limits on
their interpretation.  Such outputs are most useful when combined with rigorous
split auditing, external validation, and biochemical experiments.

%%====================================================================
\section*{Acknowledgements}
%%====================================================================

The authors thank Meta AI Research for providing the ESM-2 pretrained protein
language model and the RDKit developers for maintaining the open-source
cheminformatics toolkit.

%%====================================================================
\section*{Supporting Information}
%%====================================================================

The Supporting Information contains: model hyperparameters and feature details
(Tables~S1--S3); computational environment (Table~S4); dataset and split
statistics (Tables~S5--S7); atom features (Table~S8); the Davis ABL1
nominal-variant input audit (Table~S9); non-mutant EGFR and BRAF case summaries
(Tables~S10--S11); complete model parameter counts (Table~S12); strict
BindingDB external-validation provenance and results (Table~S13);
%%REV-START
a random-split breadth check across model families
(Table~S14); the numerical identity-stratified cold-target results
(Table~S15); matched-protocol comparison provenance (Table~S16); and the
complete BindingDB filtering waterfall for both snapshots (Table~S17).
%%REV-END
This material is available free of charge.

%%====================================================================
\section*{Data Availability}
%%====================================================================

All data used in this study are publicly available. The Davis kinase inhibitor
dataset used for model training and evaluation is available from the original
publication.\cite{davis2011comprehensive} The external evaluation uses the official
BindingDB curated-articles TSV snapshot
\texttt{BindingDB\_BindingDB\_Articles\_202607}, obtained from the BindingDB
download service;\cite{gilson2016bindingdb} the raw archive remains available from
its source, and the release provides the deterministic curation and frozen
evaluation scripts plus manifest hashes.
%%REV-START
The verification waterfall additionally uses the complete
\texttt{BindingDB\_All\_202608} release from the same download service.
%%REV-END
The immutable release containing the
released checkpoints, canonical prediction CSVs, provenance hashes, strict-split
artifacts, and reproduction scripts is archived at Zenodo,
\url{https://doi.org/10.5281/zenodo.21437753}. The BioInteract source repository
is hosted on \href{https://github.com/Liubo-520/bio-intelligence}{GitHub}. An interactive web
server for the custom-prediction workflow is available through the
\href{https://huggingface.co/spaces/AI4deeperScience/BioInteract}{BioInteract Space}; the README
documents the command-line attribution workflow.

\section*{Code Availability}

The BioInteract source code, pretrained model weights, configuration files,
canonical metric-reconstruction scripts, and figure-generation scripts are
publicly available through \href{https://github.com/Liubo-520/bio-intelligence}{GitHub} and
as the immutable Zenodo release, \url{https://doi.org/10.5281/zenodo.21437753}.
%%REV-START
The matched-protocol baseline implementations, the BindingDB filtering-waterfall
audit, and the applicability-domain analysis added in this revision are included
in the same repository and archive.
%%REV-END
The interactive web
application for the custom-prediction workflow is available through the
\href{https://huggingface.co/spaces/AI4deeperScience/BioInteract}{BioInteract Space}.

%%====================================================================
\section*{Author Contributions}
%%====================================================================

\textbf{Shengnan Wang}: Conceptualisation, Methodology, Software, Formal
analysis, Writing -- original draft. \\
\textbf{Qi Zhang}: Methodology, Validation, Data curation, Visualisation,
Writing -- original draft. \\
\textbf{Xucheng Liu}: Investigation, Validation, Data curation. \\
\textbf{Zixuan Wang}: Software, Investigation, Data curation. \\
\textbf{Jiangke Cheng}: Resources, Validation, Visualisation. \\
\textbf{Zhi Huang}: Supervision, Writing -- review \& editing, Funding
acquisition, Corresponding author.

\vspace{0.5em}
\noindent S.\,Wang and Q.\,Zhang contributed equally to this work.

%%====================================================================
\section*{Declaration of Competing Interests}
%%====================================================================

The authors declare no competing interests.

\section*{Funding}

This work was supported by the Sichuan Provincial Key Laboratory for Development and Utilization of Characteristic Horticultural Biological Resources under Grant TSYY202502. The funder had no role in study design, data collection and analysis, decision to publish, or preparation of the manuscript.

\bibliography{references}
\end{document}
